% Physical pages 3 and 4 of the concise study: the thematic orientation, in the
% order of the Mass, with what the Fathers say of each text. Every claim and
% quotation comes from the reviewed expansive study. The developed comparison
% of the three interpretations begins on the fifth page.

\runninghead{Themes and Movement}
\section{The Propers: Themes and Movement}
\zlabel{triptych:concise:themes:start}

In the Gospel of this Mass Jesus tells the crowds two short parables: a
mustard seed, ``the least indeed of all seeds,'' which ``becometh a tree,'' and
leaven that a woman ``hid in three measures of meal, until the whole was
leavened'' (Mt 13:31--33). In the Epistle St Paul gives thanks for a church at
Thessalonica that received the word ``in much tribulation'' and now waits for
the Son from heaven. Three interpretations of the Mass answer three questions
that these texts raise. The first asks how the least of seeds becomes a tree.
St Hilary of Poitiers says that the seed is the Lord, whose strength ``is
kindled by tribulations and pressures,'' and St Ambrose that mustard gives out
its strength when it is ground, as faith does when adversity presses it; on
this interpretation the Thessalonians, pressed first and then heard of in
every place, grew in the way Ambrose describes. The second asks what the leaven
and the three measures are. St Augustine and St Bede say that the leaven is
love, St Jerome that it is the knowledge of the Scriptures, Ambrose the
Church's spiritual teaching, and all four that the measures are the whole person; this
interpretation takes the charity that Paul remembers in the Thessalonians for
that leaven at work. The third asks why the chants and the Epistle speak of a
deliverance as done while those who sing them still wait for it. Augustine
says that what God has foretold is as certain as if it were past; Jerome, that
God, who could have kept his promise to the exiles without their prayer, urges
them to pray so that they may deserve to receive it.

\subsection{Introit and Collect: a letter to exiles, and a prayer for word and deed}

The Introit comes from a letter that Jeremias sent from Jerusalem to the
exiles in Babylon. It told them to build houses and plant orchards there, to
``seek the peace of the city, to which I have caused you to be carried away
captives; and pray to the Lord for it,'' and promised that ``when the seventy
years shall begin to be accomplished in Babylon'' God would bring them back
(Jer 29:5--10). The antiphon opens \latin{Dicit Dóminus}, ``The Lord saith,''
and keeps three things God says he will do: think thoughts of peace, hear
those who call on him, and bring back the captivity ``from all places.'' It
leaves out ``to give you an end and patience,'' which Jerome glosses as ``the
end of the captivity, and patience in the present labours'' (\work{In
Hieremiam}~V, PL~24 col.~893). The antiphon's verbs are future,
\latin{invocábitis \dots\ exáudiam \dots\ redúcam}; its psalm verse is past,
``Thou hast turned away the captivity of Jacob'' (Ps 84:2). Augustine takes
the verse as said to Christ, who turned away ``our captivity, not by setting
us free from the barbarians \dots\ but by setting us free from bad works, from
our sins'' (\work{En.\ in Ps.}~84.2).

The Collect has one petition. Those who pray are described as \latin{semper
rationabília meditántes}, ``ever fixing our thoughts on reasonable things''
(\latin{semper}, always; \latin{meditántes}, meditating), that is, on things
that accord with reason; they ask to do what pleases God \latin{et dictis
\dots\ et factis}, ``both in word and in deed.''

\subsection{Epistle: a church founded in a few weeks and at once persecuted}

Paul had reasoned in the synagogue of Thessalonica for three sabbaths that the
Christ had to suffer and rise; then a mob attacked the house of Jason where he
lodged, and the brethren sent him away by night to Berea (Acts 17:1--10). The lesson is his thanksgiving
to that church. He remembers \latin{óperis fidei vestræ, et labóris, et
caritátis}, ``the work of your faith and labour and charity,'' and ``the
enduring of the hope'' (1~Thess 1:3). The gospel came to them ``not \dots\ in
word only, but in power also'' (1:5); they received it ``in much tribulation,
with joy of the Holy Ghost'' (1:6); and from them ``was spread abroad the word
of the Lord \dots\ in every place'' (1:8). They turned from idols ``to serve
the living and true God. And to wait for his Son from heaven \dots\ Jesus, who
hath delivered us from the wrath to come'' (1:9--10). The last verb,
\latin{erípuit}, is a perfect, and the wrath, \latin{ira ventúra}, is still to
come.

St John Chrysostom answers the question how tribulation and joy can meet:
``The affliction is in things bodily, and the joy in things spiritual''
(\work{Hom.\ in 1~Thess.}~1). In verse~3 the Latin has three genitives side by side, work,
labour and charity; the Greek that Chrysostom expounds reads ``the labour of
love.''

\subsection{Gradual and Alleluia: thanks from a psalm of complaint, and a cry from the depths}

The Gradual sings two verses of Psalm~43: ``Thou hast delivered us, O Lord,
from them that afflict us: and hast put them to shame that hate us. In God we
will glory all the day.'' The psalm first recalls that the fathers ``got not
the possession of the land by their own sword'' (43:4), and after these verses it
complains: ``But now thou hast cast us off, and put us to shame'' (43:10);
``for thy sake we are killed all the day long: we are counted as sheep for the
slaughter'' (43:22). Paul quotes that verse to the Romans (Rom 8:36), and
because he did so Augustine hears in the psalm ``the voice of the Martyrs''
(\work{En.\ in Ps.}~43.1). The Gradual keeps the thanks and leaves the
complaint. Its \latin{Liberásti} and \latin{confudísti} are past; its
\latin{laudábimur}, ``we will glory,'' is future.

The Alleluia sings the opening of Psalm~129: ``From the depths I have cried to
Thee, O Lord: Lord, hear my prayer.'' The chant has \latin{oratiónem meam},
``my prayer,'' where the Vulgate has ``my voice.'' Hilary places the
one who cries in three depths: the depth of the riches of God, the surge of
human calamities, and the secret of the heart in which we are commanded to
pray (Mt 6:6; \work{Tractatus in Ps.}~129.2).

\subsection{Gospel: the mustard seed and the leaven}

The parables belong to the discourse of Matthew~13, which Jesus speaks from a
boat to a crowd on the shore. He tells the parable of the sower and that of
the cockle, then these two, and explains only the cockle to the disciples in
the house afterwards (Mt 13:36). Chrysostom supposes that they let the mustard
seed and the leaven pass ``as being plainer'' (\work{Hom.\ in Mt.}~47.1). Each
parable has something unexpected in it. Mustard is a garden herb, and the
parable says that when grown it is ``greater than all herbs'' before it says
\latin{et fit arbor}, ``and becometh a tree.'' Leaven elsewhere stands for
corruption: ``beware of the leaven of the Pharisees'' (Mt 16:6); ``Purge out
the old leaven'' (1~Cor 5:7). Here the kingdom of heaven is like it. The woman
\latin{abscóndit}, ``hid,'' the leaven in three \latin{sata}, a large
quantity, since Jerome notes that the \latin{satum} held a \latin{modius} and a
half (\work{In Matthaeum}~II, PL~26 col.~92); it works \latin{donec
fermentátum est totum}, ``until the whole was leavened.'' Matthew closes with a
psalm he calls the prophet's, Ps 77:2 in a form of his own: ``I will utter
things hidden from the foundation of the world.''

\subsection{Offertory to Postcommunion: the cry again, the offering, the promise and the hunger}

The Offertory sings the Alleluia's verses again and repeats the first line, so that \latin{de profúndis clamávi} is
heard three times before the Secret. The Secret asks that ``this oblation
\dots\ cleanse, renew, govern, and protect us'': two pairs of verbs, the first
asking for a change in those who offer, the second for their keeping
afterwards.

The Communion is the Lord's promise in Mark, spoken on the road to Jerusalem
in the last days before the Passion, the morning after he cursed the fig tree.
The verse before it says that whoever ``shall not stagger in his heart, but
believe'' can bid a mountain be cast into the sea (Mk 11:23); the verse after
it says, ``when you shall stand to pray, forgive, if you have aught against
any man'' (11:25). The antiphon sings only the middle verse, a command in the
present and a gift in the future: \latin{crédite quia accipiétis}, ``believe
that you shall receive.'' The Postcommunion is said by those who have been
``fed with celestial delights,'' and asks that ``we may ever hunger after those
things by which we truly live.''
\zlabel{triptych:concise:themes:end}
